\documentclass[10pt,a4paper]{article}
\usepackage{fontspec,amsmath,amssymb,graphicx}
\usepackage[margin=18mm]{geometry}
\setmainfont{Times New Roman}
\setlength{\parindent}{0pt}
\setlength{\parskip}{5pt}
\allowdisplaybreaks
\sloppy
\emergencystretch=2em
\begin{document}
\section*{BT\_3\_1\_3\_05 — Lời giải bổ sung}
\subsection*{Kiến thức và trình tự giải}
\textbf{Bước 1 — Nhận dạng dữ liệu và tốc độ.} Viết $u_{tt}-a^2u_{xx}=F$, $u(x,0)=f(x)$, $u_t(x,0)=g(x)$; tốc độ là $a>0$, không phải $a^2$. Với dữ liệu trơn, cần $f\in C^2$, $g\in C^1$. Nếu dữ liệu có góc hoặc bước nhảy, công thức vẫn mô tả sóng từng miền nhưng phải xác định nơi mất tính trơn.

\textbf{Bước 2 — Tìm hai sóng chạy.} Đặt $\xi=x-at$, $\eta=x+at$. Quy tắc dây chuyền cho $u_{tt}-a^2u_{xx}=-4a^2u_{\xi\eta}$, nên khi $F=0$, $u=A(\xi)+B(\eta)$. Tại $t=0$, $A+B=f$, $a(B'-A')=g$. Chọn một nguyên hàm $H'=g$ rồi lấy $A=f/2-H/(2a)$, $B=f/2+H/(2a)$; hằng số của $H$ sẽ triệt tiêu.


\[
u_0(x,t)=\frac{f(x-at)+f(x+at)}2+\frac{H(x+at)-H(x-at)}{2a}.
\]


\textbf{Bước 3 — Tính tích phân theo dữ liệu thật.} Tìm nguyên hàm trên từng khoảng, nối liên tục nếu $g$ có bước nhảy; thay hai cận $x\pm at$. Dùng công thức cộng lượng giác trước khi kết luận chu kỳ. Với dữ liệu chỉ thị, tích phân bằng độ dài giao của $[x-at,x+at]$ với đoạn hỗ trợ, nhân hệ số của $g$.

\textbf{Bước 4 — Nếu có nguồn, thêm Duhamel.} Nghiệm riêng có dữ liệu đầu bằng không là


\[
u_F(x,t)=\frac1{2a}\int_0^t\int_{x-a(t-s)}^{x+a(t-s)}F(y,s)\,dy\,ds.
\]


Tính tích phân trong trước. Nếu $F$ là đa thức hoặc một mode lượng giác, có thể đoán $u_F$ rồi giải ODE theo $t$, nhưng vẫn phải buộc $u_F(x,0)=(u_F)_t(x,0)=0$. Nghiệm cuối là $u_0+u_F$.

\textbf{Bước 5 — Kiểm tra và đọc ý nghĩa.} Đạo hàm công thức: $u_x=(f'(\xi)+f'(\eta))/2+(g(\eta)-g(\xi))/(2a)$, $u_t=a(f'(\eta)-f'(\xi))/2+(g(\eta)+g(\xi))/2$. Tại $t=0$ thu lại $f,g$; đạo hàm tiếp cho PDE. Các điểm gãy $s=b$ phát sinh đường $x\pm at=b$. Chu kỳ phải đúng cho mọi $x,t$, không chỉ sau khi một mặt sóng đi qua một điểm cố định.
\subsection*{Lời giải áp dụng cho bài}
\textbf{Bước 1. Giải phần thuần nhất.} Thay $f,g$ vào d'Alembert. Nguyên hàm cần dùng lần lượt là $-\cos2s/2$, $\arctan s$, $-\cos s$, $e^s$, $\sin s$ và $H$ của xung tam giác. Sau đó tìm phần nguồn có hai dữ liệu đầu bằng không.

\textbf{Bước 2. Tính từng nguồn.}

(i) $a=2$, nguồn $x\cos t$ có nghiệm riêng $x(1-\cos t)$ vì đạo hàm không gian cấp hai bằng không. Suy
\[u=e^{-x}\cosh2t+\tfrac14\sin2x\sin4t+x(1-\cos t).\]
(ii) $a=3$. Nguồn $e^x$ dùng $e^xq$, $q''-9q=1$, $q(0)=q'(0)=0$, cho $q=(\cosh3t-1)/9$. Nguồn $\sin t$ độc lập $x$ cho $t-\sin t$:
\[u=\frac{\arctan(x+3t)-\arctan(x-3t)}6+\frac{e^x(\cosh3t-1)}9+t-\sin t.\]
(iii) $a=4$, nguồn $x+t$ cho $xt^2/2+t^3/6$:
\[u=x^2+16t^2+\tfrac14\sin x\sin4t+xt^2/2+t^3/6.\]
(iv) $a=7$, phần vận tốc $e^x$ là $e^x\sinh7t/7$, phần nguồn $xt$ là $xt^3/6$.

(v) $a=1$, $F=\sin(x+t)=\sin x\cos t+\cos x\sin t$. Hai ODE là $q''+q=\cos t$ và $r''+r=\sin t$ với dữ liệu không, ở cộng hưởng. Suy $q=t\sin t/2$, $r=(\sin t-t\cos t)/2$:
\[u=x+\cos x\sin t+\tfrac12\{\sin x\,t\sin t+\cos x(\sin t-t\cos t)\}.\]
(vi) $a=2$, $f$ là hình thang trên $[0,3]$ với nhánh $s,1,3-s$ tại các đoạn $[0,1],[1,2],[2,3]$. $H$ là nguyên hàm tam giác ở Bài 2(vi). Nguồn $-e^x$ cho
\[u=\tfrac12[f(x-2t)+f(x+2t)]+\frac{H(x+2t)-H(x-2t)}4-\frac{e^x(\cosh2t-1)}4.\]
\textbf{Bước 3. Kiểm tra.} Mỗi phần nguồn và đạo hàm thời gian của nó bằng không ở $t=0$. Các ODE vừa lập chứng minh PDE, không chỉ đoán dạng đáp án. Câu (vi) có góc của $f,g$, nên nghiệm được hiểu từng miền/ theo nghĩa yếu trên các đường $x\pm2t=0,1,2,3$ và $x\pm2t=-1,0,1$.

\textbf{Bước 4. Vẽ đặc trưng và các lát cắt.} Với mỗi tốc độ $a$ của từng câu, các đường xuất phát từ $(b,0)$ là $x=b+at$ và $x=b-at$. Nguồn thay đổi giá trị trên miền ảnh hưởng nhưng không đổi hai họ đặc trưng này. Hình kèm lấy $t=0,0.15,0.35,0.65$ và thay trực tiếp vào sáu công thức đã tính; hình thứ hai cho mặt phẳng $x,t$ với các đường đặc trưng mẫu.

\textbf{Chi tiết giải các ODE nguồn.} (i) Đặt $u_F=xq$, có $q''=\cos t$, $q'(0)=q(0)=0$. Tích phân $q'=\sin t$, $q=1-\cos t$. (ii) Với $u_F=e^xq$, $q''-9q=1$: thử $q_p=-1/9$, phần thuần nhất $A\cosh3t+B\sinh3t$. Dữ liệu cho $A=1/9,B=0$. Nguồn $\sin t$ cho $r''=\sin t$, $r'=1-\cos t$, $r=t-\sin t$.

(iii) $q''=x+t$, tích phân theo thời gian với $x$ cố định được $q'=xt+t^2/2$, $q=xt^2/2+t^3/6$. (iv) $q''=xt$ cho $q'=xt^2/2$, $q=xt^3/6$.

(v) Tính trực tiếp $q=t\sin t/2$: $q'=(\sin t+t\cos t)/2$, $q''=\cos t-t\sin t/2$, nên $q''+q=\cos t$. Với $r=(\sin t-t\cos t)/2$, $r'=t\sin t/2$, $r''=(\sin t+t\cos t)/2$, nên $r''+r=\sin t$. Cả $q,r$ và đạo hàm đầu bằng $0$. (vi) $u_F=e^xq$ cho $q''-4q=-1$, $q_p=1/4$, $q=A\cosh2t+B\sinh2t+1/4$; dữ liệu cho $A=-1/4,B=0$.
\begin{center}\includegraphics[width=.97\linewidth,height=11cm,keepaspectratio]{Do_thi/BT_3_1_3_05.png}\end{center}
\textit{Đồ thị minh họa; với nghiệm số, xem phương pháp và tham số trong mã Python kèm theo.}
\begin{center}\includegraphics[width=.97\linewidth,height=11cm,keepaspectratio]{Do_thi/BT_3_1_3_05_characteristics.png}\end{center}\end{document}